\section{Infinite paths and weak-Cauchy block sequences}
\label{sec:paths}

An infinite path is an infinite list of pieces, declared weights and
metadata whose every finite prefix satisfies the rules of
Section~\ref{sec:norming}.  We first show that different paths have
almost disjoint weight sets.  We then use the Type~II operation to obtain
convergence on all paths simultaneously, and show that this already
controls every functional on a block subspace.

\begin{lemma}[Bounded path functionals]\label{path:bounded}
For every actual infinite path $(\phi_l,j_l,p_l)_{l\ge1}$, the formula
\[
 P(x)=\sum_{l=1}^{\infty}\phi_l(x)
\]
defines a member of $B_{E^*}$.  The sum converges for every $x\in E$.
The same is true after an allowed crop of the path.
\end{lemma}

\begin{proof}
Every finite partial sum is a Type~II functional, with one path and
coefficient one.  These sums have norm at most one.  Their evaluations
eventually stabilize on every finitely supported vector, since the
pieces have successive disjoint supports.  Uniform boundedness of the
partial sums and finite approximation therefore give convergence on
every vector in $E$, with limit norm at most one.  Cropping follows
either by the same argument or by composing with a contractive
projection.
\end{proof}

\begin{lemma}[Rigidity of shared weights]\label{path:rigidity}
Distinct infinite paths in the same family have only finitely many
weights in common.  Consequently, for any finitely many distinct paths
$P_1,\ldots,P_h$ in that family, one may delete finitely many initial
pieces of each to obtain tails $\widetilde P_1,\ldots,
\widetilde P_h$ satisfying
\begin{equation}
 \left\|\sum_{b=1}^{h}a_b\widetilde P_b\right\|_{E^*}
 \le\|(a_b)_{b=1}^{h}\|_{\ell_q}
 \qquad(a_b\in\mathbb R).
 \label{path:tail-estimate}
\end{equation}
Each difference $P_b-\widetilde P_b$ belongs to $X_0$.
\end{lemma}

\begin{proof}
Suppose two infinite paths have infinitely many common weights.
Except possibly for their first weights, equality of weights implies
equality of the preceding codes, by injectivity of $\sigma$.
In particular, the preceding prefixes have the same length, the same
metadata and weights, and the same arrays in their respective ordered
closures.  The lengths of these matching prefixes tend to infinity.

Interchange the two paths if necessary so that the supremum of the
support of the first does not exceed that of the second.  Fix a
coordinate $\alpha$ in the first path.  There is a coordinate
$\xi\ge\alpha$ in the second path: an infinite successive path has no
greatest support coordinate.  For sufficiently long matching prefixes,
both $\alpha$ and $\xi$ have appeared, and their common final metadata
$p$ is at least $\rho(\alpha,\xi)$.  The closure for the second prefix
contains $\xi$, hence $F_p(\xi)$, hence $\alpha$; the closure for the
first prefix also contains $\alpha$.  Their ordered initial segments
through $\alpha$ coincide by Lemma~\ref{ch:closures}.

Every piece preceding the first path's piece that contains $\alpha$
is supported strictly before $\alpha$.  Position alignment therefore
makes this entire piece exactly equal to the corresponding piece in
the second path, not merely equal on an initial interval.  Its weight
and metadata agree as well.  Taking $\alpha$ in arbitrarily late
pieces identifies all the path data.  Thus the paths were identical.

For finitely many distinct paths, discard initial pieces beyond all
pairwise common weights.  Their remaining weight sets are disjoint.
Each finite sum of tail pieces is the crop of a finite prefix of the
original path; this retains the required uncropped coding.  Type~II
therefore gives \eqref{path:tail-estimate} first for rational coefficients
and finite tails.  Pass to the path limits on each vector, use the norm
bound, and approximate real coefficients.  The discarded pieces have
finite support and hence belong to $X_0$.
\end{proof}

Fix a pure family $k$ or the mixed family.  By a block sequence in the
family we mean a sequence of nonzero finitely supported vectors that
are successive and supported on $C_k$ in the pure case, or
mixed-successive in the mixed case.  A sequence $(x_i)$ is weak-Cauchy
if $(f(x_i))$ converges for every continuous linear functional $f$.

\begin{proposition}[Weak-Cauchy blocks]
\label{path:weak-cauchy}\label{path:cauchy}
Every bounded block sequence in a fixed family has a weak-Cauchy
subsequence.
\end{proposition}

\begin{proof}
Write $M=\sup_i\|x_i\|$.  The proof has two steps.  First we pass to a
subsequence on which every actual infinite path in the family
converges.  Then we show that every pointwise limit of norming
functionals converges on the same sequence.

\smallskip\noindent
\emph{Step 1: convergence on all actual paths.}
Fix $\delta>0$.  If an as yet untreated path satisfies
$\limsup_i|P(x_i)|>\delta$, pass to an infinite subsequence on which
$|P(x_i)|>\delta$ and $P(x_i)$ converges.  Continue this procedure,
retaining all previously obtained convergence.

Only finitely many new paths can be selected at this threshold.
Indeed, if $h$ distinct paths have been selected, remove their finite
prefixes as in Lemma~\ref{path:rigidity}.  Those prefixes vanish on all
sufficiently late blocks.  By \eqref{path:tail-estimate} and finite
dimensional $\ell_q$ duality, at each such block
\[
 \left(\sum_{b=1}^{h}|P_b(x_i)|^{q'}\right)^{1/q'}
 \le\|x_i\|\le M.
\]
Since every selected absolute value exceeds $\delta$ on the retained
subsequence, $h\le(M/\delta)^{q'}$.  The selection must therefore stop.
At that point every unselected path has limsup of its absolute
evaluations at most $\delta$.

Apply this finite procedure successively to thresholds decreasing to
zero and nested subsequences.  Take a diagonal subsequence with the
usual tail containment in each preceding subsequence.  Every selected
path converges on it.  An unselected path has absolute limsup at most
every threshold and hence tends to zero.  Relabel this subsequence as
$(x_i)$.

\smallskip\noindent
\emph{Step 2: convergence on pointwise limits of norming functionals.}
Each individual norming functional is finitely supported and already
vanishes on all sufficiently late blocks. This does not control every
functional on the block span. We must instead control the compact
pointwise closure of the norming set; at the end of the proof a measure
representation will pass from this closure to the whole dual.

Let $S=\bigcup_i\operatorname{supp}x_i$, a countable set.  Use $K_k$ in
the pure case and $K$ in the mixed case, and let $\mathcal C$ be the
pointwise closure of its restrictions in $[-1,1]^S$.  This is a compact
metrizable space.  Each $u\in\mathcal C$ acts with norm at most one on
the closed block span, because this holds on finite vectors by
pointwise approximation.  Define its oscillation on the blocks by
\[
 \operatorname{osc}(u)
 =\limsup_i u(x_i)-\liminf_i u(x_i),\qquad
 D=\sup_{u\in\mathcal C}\operatorname{osc}(u)\le2M.
\]
We will prove that $D=0$.

Fix $u\in\mathcal C$ and choose norming functionals converging
pointwise on $S$ to $u$.  Subsequences may be taken without changing
this limit.  Choose one valid expression for each functional.  After
subselection the top expressions are all terminals, all Type~I, or
all Type~II.  The following analysis does not require their
construction depths to be bounded: all child limits will lie in the
same compact set $\mathcal C$.

\smallskip\noindent
\emph{Terminals and Type~I limits.}
In a pure family, a sequence of coordinate terminals has either a
fixed coordinate on a subsequence or the zero pointwise limit; its
oscillation is zero.  In the mixed family, terminal functionals belong
to individual pure colors.  If those colors drift, the pointwise
limit is zero; if the color is fixed, it meets at most one of the
mixed-successive blocks, again giving zero oscillation.

For Type~I expressions whose weights tend to infinity, the coordinate
bound $1/m_j$ gives pointwise limit zero.  Otherwise fix the weight
$j$, the number of children, and the pointwise limits of all children.
Thus
\[
 u=\frac1{m_j}\sum_{h=1}^{d}v_h,
 \qquad d\le L_j,\quad v_h\in\mathcal C.
\]
Nonzero coordinates of different child limits remain in their
original order.  At most one child limit can meet infinitely many
blocks.  For if an earlier child met infinitely many, its hits would
eventually occur beyond any fixed hit of a later child, contradicting
the order.  Every other child is zero on all sufficiently late
blocks.  Consequently
\begin{equation}
 \operatorname{osc}(u)\le D/m_j.
 \label{path:type-i-oscillation}
\end{equation}
This conclusion also applies to a pointwise limit of cropped
Type~I pieces: distribute the crops to the children and discard zeros
before taking the limits.

\smallskip\noindent
\emph{Individual terms in a Type~II limit.}
The key issue is the limit of one cropped path. We shall show that, on
all sufficiently late blocks, such a limit agrees either with zero,
with a limit of a single Type~I piece, or with an actual infinite path.
The first and third alternatives have zero oscillation; the second has
the small factor $1/m_j$ from \eqref{path:type-i-oscillation}. After
classifying the individual terms, we will use their disjoint active
weights to sum these factors.

Index each cropped path term by its smallest active weight $b$.
These indices are distinct within each expression.  A diagonal
subsequence makes presence or absence eventually constant for each index,
makes its coefficient converge to a number $a_b$, and makes its crop
membership converge at every coordinate of $S$.  Put $a_b=0$ for
absent terms.  Fatou's inequality on finite sets gives
\begin{equation}
 \sum_b|a_b|^q\le1.
 \label{path:limit-coefficients}
\end{equation}
For each present term, also take pointwise limits of its path crop.
Call the resulting functional on the block span $u_b$.
All subsequent choices impose only countably many conditions, indexed
by a term, a finite chain position and a coordinate of $S$. Nested
subsequences and a diagonal retain every finite collection of these
conditions simultaneously for all sufficiently late approximants.

The original uncropped weight list through $b$ has only finitely many
possibilities, since weights strictly increase.  Fix this initial
list by subselection.  Examine its successive extensions one position
at a time.  At each position the path either ends, the next weight
tends to infinity, or the next weight can be fixed on a subsequence.
Diagonalizing produces a finite fixed chain, possibly followed by a
divergent-weight suffix, or an infinite fixed chain.  Take pointwise
limits of all cropped pieces in this chain as well.

A suffix whose weights tend to infinity contributes zero in the
pointwise limit.  Each coordinate meets at most one piece of the
suffix, and its coefficient is bounded by the reciprocal of its
least weight's $m$-parameter.  Whenever a fixed piece has a fixed
successor weight, that successor fixes its entire preceding code.
In particular, the support size of the piece is uniformly bounded
by the length of the finite array in that code.  Its pointwise limit
therefore has finite support.

For a finite fixed chain, only its last piece can have a limit with
infinite support.  If its weight is $j$,
\eqref{path:type-i-oscillation} gives
\begin{equation}
 \operatorname{osc}(u_b)\le D/m_j.
 \label{path:finite-chain-oscillation}
\end{equation}
If the last cropped piece has zero pointwise limit, then $u_b$ is
finitely supported and its oscillation is zero.  Otherwise the last
weight is active in all sufficiently late approximating terms.
The last weights for these potentially oscillating terms are
therefore distinct: two of them cannot be active in different terms
of the same Type~II approximant.

\smallskip\noindent
\emph{Infinite fixed chains and their limiting crops.}
It remains to transfer convergence from actual paths to infinite fixed
chains, then sum the Type~II terms in dual norm.
In this case every individual piece has a finitely supported
pointwise limit.  The pointwise limit of the path crop is the sum of
these piece limits.  To see that no additional contribution escapes
to high piece positions, note that the coordinate error after the
first $l$ pieces is at most $1/m_{j_{l+1}}$, uniformly in sufficiently
late approximants.  If only finitely many cropped piece limits are
nonzero, $u_b$ is finitely supported and has zero oscillation.

Suppose instead that arbitrarily late pieces have nonzero cropped
limits.  Choose a coordinate $\beta\in S$ with nonzero limit in one
such piece.  Eventually $\beta$ belongs to the actual support of
that piece and is retained by its crop.  The fixed next weight fixes
the prefix code through this piece, including its metadata $p$.
All corresponding approximating closures contain $\beta$, and their
initial segments through $\beta$ are therefore exactly
$F_p(\beta)$.  Every earlier uncropped piece is supported strictly
before $\beta$.  The position-alignment observation following
Lemma~\ref{ch:closures} shows that these earlier pieces are eventually
the same actual functions on $\Gamma$.

Choosing $\beta$ in arbitrarily late effective pieces proves eventual
stabilization of every uncropped prefix, with its weights and
metadata.  These stabilized prefixes define an actual infinite path
$P$.  Its pieces are valid Type~I expressions because each stabilized
piece equals a piece from an approximating path; every finite prefix
satisfies the required coding rule for the same reason.

The limiting crop changes $P$ on only finitely many coordinates of
$S$.  To prove this, fix one retained coordinate $\beta$ in a piece.
For a coordinate $\alpha\in S$ belonging to any strictly later
piece, choose a retained coordinate $\gamma$ in a still later
effective piece.  In late approximants the stabilized coordinates
$\beta,\alpha,\gamma$ have their fixed ordinal order, and, in the
mixed case, their fixed color order.  Both $\beta$ and $\gamma$ lie
in the crop.  Interval convexity, in both orders when required,
forces $\alpha$ to lie in it too.  Thus the limiting crop retains
every path coordinate of $S$ after the piece containing $\beta$.
Only the finitely many finite pieces up to and including the piece
containing $\beta$ can differ.
It follows that $u_b(x_i)-P(x_i)$ is eventually zero.  Step~1 proves
convergence of $P(x_i)$, so every infinite-chain term has zero
oscillation.

\smallskip\noindent
\emph{Reassembling the Type~II limit.}
We have
\begin{equation}
 u=\sum_b a_bu_b
 \quad\hbox{pointwise on }S.
 \label{path:limit-series}
\end{equation}
Indeed, terms with smallest active weight $b>N$ have coordinate
error at most
\[
 \sum_{j>N}\frac1{m_j},
\]
uniformly in all approximants.  At a given coordinate there is at
most one contributing piece per path, and the contributing active
weights are globally distinct.  This proves
\eqref{path:limit-series} after passing to the finitely many indices
$b\le N$ and then letting $N\to\infty$.

The series also converges in dual norm on the block span.  For every
finite set of present indices and every real coefficient vector
$(c_b)$ with $\ell_q$ norm at most one, replace the coefficients of
the corresponding approximating Type~II terms by rational
approximations to $(c_b)$ inside the closed $\ell_q$ unit ball.
They remain admissible, since the paths
and their pairwise disjoint active weights have not changed.
Passing to the pointwise limits on finite vectors yields
\[
 \left\|\sum_b c_bu_b\right\|_{[x_i]^*}\le1.
\]
By \eqref{path:limit-coefficients}, finite truncations of
$\sum_ba_bu_b$ are therefore Cauchy in dual norm.  Their norm limit
agrees with \eqref{path:limit-series} on the block vectors.

Oscillation is subadditive and changes by at most
$2M\|v-w\|$ when one bounded functional is replaced by another.
Using finite truncations, the distinct potentially oscillating last
weights in \eqref{path:finite-chain-oscillation}, and $|a_b|\le1$,
we obtain
\begin{equation}
 \operatorname{osc}(u)\le D\sum_{j=1}^{\infty}\frac1{m_j}.
 \label{path:type-ii-oscillation}
\end{equation}
All other terms have zero oscillation.

The terminal, Type~I and Type~II cases are exhaustive.  Equations
\eqref{path:type-i-oscillation} and
\eqref{path:type-ii-oscillation}, together with
Lemma~\ref{ch:parameters}, imply
\[
 D\le\left(\sum_j1/m_j\right)D<\tfrac18D
 \quad\hbox{if }D>0.
\]
Hence $D=0$: every $u\in\mathcal C$ gives a convergent sequence
$u(x_i)$.

Finally, evaluation embeds the block span isometrically into
$C(\mathcal C)$; for a finite vector, its evaluation is continuous,
and finite approximation extends this assertion to the closed span.
Given any functional on the block span, extend it by Hahn--Banach to
$C(\mathcal C)$. The Riesz--Markov representation theorem
\cite[Theorem~6.19]{Rudin1987} represents the extension by a finite
signed Borel measure $\nu$, with $|\nu|(\mathcal C)<\infty$.
The functions $u\mapsto u(x_i)$ are uniformly bounded by $M$ and
converge at every $u\in\mathcal C$. Bounded convergence with respect
to $|\nu|$ proves convergence under the given functional.  The selected subsequence is therefore weak-Cauchy.
\end{proof}
